\documentclass[10pt]{article}
\usepackage[T1]{fontenc}
\usepackage{lmodern}
\usepackage[margin=0.82in]{geometry}
\usepackage{amsmath,amssymb,amsthm}
\usepackage[hidelinks]{hyperref}
\setlength{\parindent}{0pt}
\setlength{\parskip}{4pt}
\newtheorem{theorem}{Theorem}
\newcommand{\Z}{\mathbb Z}
\newcommand{\Wh}{\operatorname{Wh}}
\newcommand{\GL}{\operatorname{GL}}
\begin{document}
\begin{center}
{\Large A finite hyperbolic group with nonzero Whitehead group}\\[4pt]
{\large Disproof of the unrestricted vanishing clause in TLMC 00000001520}\\[3pt]
\texttt{lizaixi01}\qquad 24 October 2026
\end{center}

\textbf{Exact scope.}
The English source says: ``For Gromov-hyperbolic groups $\Gamma$,
$\Wh(\Gamma)=0$ (torsion-free and trivial).'' The Chinese likewise places
its parenthesis after the assertion $\Wh(\Gamma)=0$. Neither version
imposes torsion-freeness on $\Gamma$. We refute that printed unrestricted
vanishing assertion. The counterexample has torsion, so it does not
refute a corrected assertion restricted to torsion-free groups. We make
no claim about the subsequent boundary-action, $2$-torsion, or
Farrell--Jones clauses. The complete bilingual source is in \texttt{SOURCE.md}.

\begin{theorem}
Let $G=C_5=\langle t\mid t^5=1\rangle$ and $A=\Z[G]$. Then $G$ is
a finitely generated Gromov-hyperbolic group and $\Wh(G)\ne0$.
More precisely, the genuine group-ring unit
\[
 u=1-t-t^{-1},\qquad v=1-t^2-t^{-2}
\]
satisfies $uv=vu=1$ and has nonidentity image in $\Wh(G)$.
\end{theorem}

\textbf{1. The actual stable groups and Whitehead quotient.}
For an arbitrary commutative ring $R$, take finite ordered index sets
$I_0=\{*\}$, $I_{n+1}=I_n\sqcup\{*\}$. Their proved cardinalities are
$n+1$. The maps
\[
 \GL(I_n,R)\longrightarrow\GL(I_{n+1},R),\qquad
 M\longmapsto\begin{pmatrix}M&0\\0&1\end{pmatrix}
\]
are actual injective homomorphisms. Lean constructs the whole
$\mathbb N$-indexed functor and its native filtered colimit $\GL_\infty(R)$
in the category of groups. We use the standard definition
$K_1(R)=\GL_\infty(R)^{\rm ab}$, realized by native abelianization.
This is the stable-abelianization definition of algebraic $K_1$;
the equivalence with the elementary-matrix definition is Whitehead's
lemma, recalled in Milnor~\cite[pp.~358--360]{Milnor}.

For $A=\Z[C_5]$, let $T\le A^\times$ be the range of the true homomorphism
\[
 \Z^\times\times C_5\longrightarrow A^\times,\qquad (\epsilon,g)\mapsto\epsilon g.
\]
The classification $\Z^\times=\{1,-1\}$ is used in Lean to prove that
$T$ is exactly the subgroup generated by $-1$ and the group elements.
If $j:A^\times\to K_1(A)$ is the actual one-by-one-matrix inclusion,
the formal Whitehead group is
\[
 \Wh(C_5)=K_1(A)/j(T).
\]
Lean also proves that $j(T)$ is the subgroup generated by the images
of $-1$ and all group elements. This is the signed quotient of the
standard definition in Milnor~\cite[p.~373]{Milnor}.

\textbf{2. A determinant detector and a concrete nonzero class.}
The ordinary determinants commute with every block inclusion. Their
cocone gives, by the native group-colimit universal property,
$\det_\infty:\GL_\infty(A)\to A^\times$. Native abelianization and
quotient constructions then give genuine homomorphisms
\[
 K_1(A)\longrightarrow A^\times,
 \qquad \Wh(C_5)\longrightarrow A^\times/T.
\]
The first map sends $j(a)$ to $a$, and the second sends the Whitehead
class of a unit $a$ to its actual coset $aT$. The unit quotient is
only a detection target; it is not substituted for $K_1$ or $\Wh$.

\newpage
In the actual group algebra, $t^{-1}=t^4$ and $t^{-2}=t^3$. Hence
\[
 (1-t-t^4)(1-t^2-t^3)
 =1+(t+t^2)(t^5-1)=1.
\]
This proves invertibility with the displayed inverse. The coefficients
of $u$ at the two distinct group elements $t$ and $t^{-1}$ are both
$-1$. Every signed group unit $\epsilon g$ is supported at one group
element. Consequently $u\notin T$. Its coset in $A^\times/T$ is
nonidentity, so the determinant detector proves that its actual
Whitehead class is nonidentity as well.

Milnor's Example~6.6 uses $t+t^{-1}-1$; our $u$ is its negative and
therefore has the same signed Whitehead class. The argument above
needs neither the classification $\Wh(C_5)\cong\Z$ nor a claim that
$SK_1$ vanishes. In particular, it does not use Milnor's general
finite-abelian-group statement in Example~1.6, whose added-in-proof
warning must be respected.

\textbf{3. The genuine Cayley metric and hyperbolicity.}
Use the finite symmetric generating set $S=C_5\setminus\{1\}$.
Lean proves $|C_5|=5$, $|S|=4$, and $\langle S\rangle=C_5$ from the
actual finite group. Its native Cayley graph is the complete graph
on these five vertices and is connected. Define $d$ as the least
length of a Cayley walk, using \texttt{SimpleGraph.dist}. Existence
of a shortest walk is proved before the distance is used, so the
disconnected-distance default is never invoked. One obtains
\[
 d(g,h)=\begin{cases}0&g=h,\\1&g\ne h.\end{cases}
\]
Lean equips the group with the resulting real-valued native
\texttt{MetricSpace}, proves left invariance, and verifies the standard
four-point inequality, for all actual group elements,
\[
 d(a,c)+d(b,d)\le
 \max\{d(a,b)+d(c,d),d(a,d)+d(b,c)\}+2.
\]
Indeed the left side is at most $2$, while the maximum is nonnegative.
This is the four-point formulation of hyperbolicity for the word
metric. The standard relationship with the geodesic Cayley-graph
formulation is described in Bowditch~\cite[\S\S2.3,6.5,6.10]{Bowditch};
his \S6.10 explicitly includes every finite group. The development
proves the full Cayley-walk metric statement directly, rather than
attaching a hyperbolicity label to an unrelated group. It does not
formalize the general equivalence of four-point and thin-triangle
definitions. The proved element $t\ne1$ with $t^5=1$ records exactly
why this counterexample is outside a torsion-free-group revision.

\textbf{Formalization and reproduction.}
The final theorem is \texttt{TLMC1520.complete\_counterexample}.
The project pins Lean~4.33.0 and Mathlib commit
\begin{center}
{\footnotesize\texttt{db584cd6d46c92f209a44c0f1c829460d327499d}}.
\end{center}
From \texttt{lean/}, run these four commands separately:
\begin{verbatim}
lean --version
lake build +Main
lake env lean -DwarningAsError=true Check.lean
lake env leanchecker --fresh Main
\end{verbatim}
The audit emits ordinary \texttt{\#print axioms} for every actual owned
declaration, including generated names. Only \texttt{propext},
\texttt{Classical.choice} and \texttt{Quot.sound} are permitted.
Source-bound logs, execution receipts, actual thread samples and a
SHA-256 manifest accompany the submission. Production timings are
diagnostic and give no controlled efficiency conclusion.

\begin{thebibliography}{2}\small
\bibitem{Milnor} J.~Milnor, \emph{Whitehead torsion},
Bulletin of the American Mathematical Society~72 (1966), 358--426,
\S1, \S6, Example~6.6.
\href{https://doi.org/10.1090/S0002-9904-1966-11484-2}{DOI: 10.1090/S0002-9904-1966-11484-2}.
\bibitem{Bowditch} B.~H.~Bowditch, \emph{A course on geometric group theory},
author notes, October~2005, \S2.3 (p.~12), \S6.5 (pp.~51--53),
\S6.10 (p.~59).
\href{https://www.bhbowditch.com/papers/bhb-ggtcourse.pdf}{Author PDF}.
\end{thebibliography}
\end{document}
